\documentclass[10pt,a4paper]{article}
\usepackage[utf8]{inputenc}
\usepackage[T1]{fontenc}
\usepackage[spanish]{babel}
\usepackage{lmodern}
\usepackage{geometry}
\geometry{margin=1.6cm}
\usepackage{booktabs}
\usepackage{tabularx}
\usepackage{enumitem}
\setlist{nosep,leftmargin=*}
\usepackage{hyperref}
\newcommand{\ceddie}{\emph{verificar con CEDDIE / política institucional vigente}}
\newcommand{\clave}{\texttt{CADI-EIC-IA-AGENTICA-01}}
\pagestyle{empty}
\begin{document}
\begin{center}
{\Large\bfseries Ficha de autorización --- Líder EIC}\\[0.3em]
{\large IA agéntica para la creación de contenido académico}
\end{center}
\vspace{-0.5em}

\noindent\textbf{CADi:} IA agéntica para la creación de contenido académico\\
\textbf{Escuela:} Escuela de Ingeniería y Ciencias (EIC) · Tecnológico de Monterrey\\
\textbf{Clave sugerida:} \clave\ \textit{(asignar clave oficial con CEDDIE / EIC)}\\
\textbf{Duración:} 40\,h totales --- \textbf{24\,h aula} (8 $\times$ 3\,h) + \textbf{16\,h fuera de aula}\\
\textbf{Modalidad:} virtual (Teams) / híbrida / presencial \textit{(según edición)} · \textbf{Cupo:} 20--30\\
\textbf{Audiencia:} profesorado de profesional y posgrado de la \textbf{EIC}\\
\textbf{Instructor:} Juliho Castillo · Correo: \underline{\hspace{4cm}} · Campus: \underline{\hspace{3cm}}

\subsection*{Por qué (STEM)}
\begin{itemize}
\item El profesorado EIC produce textbooks/ebooks, guías de lab, decks y artículos con alta exigencia de \textbf{notación, unidades y citas}.
\item El uso informal de chat suele ser fragmentario y difícil de auditar.
\item La \textbf{IA agéntica} (roles, herramientas, HITL) acelera idea $\rightarrow$ outline $\rightarrow$ borrador $\rightarrow$ revisión.
\item Formato 40\,h permite práctica real y avance de un \textbf{microartefacto} entre sesiones.
\end{itemize}

\subsection*{Qué se obtiene}
\begin{itemize}
\item \textbf{RA top:} distinguir chat vs.\ agentes; diseñar brief/orquestación; producir outline, secciones, deck o sección IMRyD con verificación; cerrar con agent log + checklist.
\item \textbf{Microartefacto EIC:} capítulo STEM corto / deck técnico / sección IMRyD / guía de laboratorio breve + agent log + checklist.
\end{itemize}

\subsection*{8 sesiones (aula)}
\begin{tabularx}{\textwidth}{@{}cXl@{}}
\toprule
\# & Sesión (3\,h) & Producto \\
\midrule
1 & Fundamentos agénticos e integridad EIC & E1 mapa + nota integridad + tema \\
2 & Briefs, prompts y orquestación & E2 brief + prompt + esquema \\
3 & Libros I --- arquitectura y outline & E3 arquitectura + outline \\
4 & Libros II --- secciones y consistencia & E4 secciones + consistencia \\
5 & Presentaciones I --- narrativa / storyboard & E5 storyboard \\
6 & Presentaciones II --- deck y notas & E6 deck MV + notas \\
7 & Artículos I --- estructura y borrador & E7 estructura / borrador \\
8 & Artículos II --- revisión y cierre & E8 microartefacto + log \\
\bottomrule
\end{tabularx}

\subsection*{Evidencia de acreditación (sugerida)}
Participación (E1--E7) + evidencia de las \textbf{16\,h} independientes + plantillas + \textbf{microartefacto final} + agent log + checklist. Umbral CEDDIE (\ceddie).

\begin{quote}
\textbf{Verificación CEDDIE:} políticas de integridad académica, uso de IA, datos y herramientas homologadas --- confirmar vigencia y enlaces oficiales \textbf{antes de cada edición}. No sustituye dictamen institucional.
\end{quote}

\subsection*{Solicitud de autorización}
\begin{tabularx}{\textwidth}{@{}lX@{}}
\toprule
Campo & Espacio \\
\midrule
Nombre del/de la líder EIC & \hrulefill \\
Cargo & \hrulefill \\
Fecha & \hrulefill \\
Firma / visto bueno & \hrulefill \\
Comentarios & \hrulefill \\
\bottomrule
\end{tabularx}

\end{document}
